\chapter{The Fundamental Group of Manifolds of Negative Curvature}
\label{chap:dc-fundamental-group-negative-curvature}

Faithful transcription of do Carmo, \emph{Riemannian Geometry}. Each numbered
statement keeps do Carmo's number, kind and (name); proofs keep his explicit
cross-references ("by Cartan's Theorem", "by \cref{lem:dc-ch12-3-1} (ii)", "Chapter 7") so the
dependency graph is recoverable from the text.

\section{Introduction}

Let $M^n$ be a complete Riemannian manifold with sectional curvature $K<0$. A
fundamental fact about the topology of $M$ is that the universal covering of $M$
is diffeomorphic to $\mathbb{R}^n$ (cf. Hadamard's Theorem, Chapter 7). In this
chapter, we obtain information about the fundamental group $\pi_1(M)$ of $M$.

The importance of studying $\pi_1(M)$, when $M$ has negative curvature, comes
from the fact that, in a certain sense, all of the topology of $M$ is contained
in $\pi_1(M)$. More precisely, in Algebraic Topology one studies certain
topological invariants, called the homotopy groups of dimension $k\geq 1$, which
generalize the fundamental group ($=$ homotopy group of dimension one). Such
groups can be defined, roughly, as homotopy classes of maps $f:S^k\to M$ of
spheres $S^k$ of dimension $k$ into $M$. It is possible to prove (see M.
Greenberg [Gb], p. 32) that if $M$ is covered by $\mathbb{R}^n$, then every such
$f$ is homotopic to a constant if $k\geq 2$. This means that the homotopy groups
of higher dimension are trivial and that, therefore, at the level of homotopy,
the information about the topology of $M$ is contained in $\pi_1(M)$.

The object of this chapter is to prove the Theorem of Preissman which states the
following: If $M$ is compact and $K<0$ then every non-trivial abelian subgroup of
$\pi_1(M)$ is infinite cyclic. This shows, for example, that the torus
$S^1\times S^1\times S^1$ whose fundamental group is
$\mathbb{Z}\oplus\mathbb{Z}\oplus\mathbb{Z}$ cannot support a metric of strictly
negative curvature.

In Section 2 we introduce some important notions and prove a theorem of E. Cartan
on the existence of closed geodesics which does not use any hypothesis on the
curvature. In Section 3 we prove the Theorem of Preissman. We prove also that if
$M$ is compact and $K<0$, then $\pi_1(M)$ is not abelian. Finally, we show that,
in the statement of Preissman's Theorem, it is possible to replace the condition
"abelian subgroup" by the weaker condition "solvable subgroup".

\section{Existence of closed geodesics}

We denote by $\pi:\tilde{M}\to M$ the universal covering of a complete Riemannian
manifold $M$ with the covering metric and denote by $A(\tilde{M})$ the group of
covering transformations of $\tilde{M}$ (covering automorphisms). Observe that
the elements in $A(\tilde{M})$, different from the identity, are isometries of
$\tilde{M}$, which have no fixed points, and that $A(\tilde{M})$ is isomorphic to
$\pi_1(M;p)$, $p\in M$ (cf. Massey [Ma]). This isomorphism depends on the choice
of the point $\tilde{p}\in\tilde{M}$, with $\pi(\tilde{p})=p$, and associates to
each $g\in\pi_1(M;p)$ an isometry $\alpha_{\tilde{p}}\in A(\tilde{M})$ defined in
the following way: If $\tilde{q}\in\tilde{M}$, join $\tilde{q}$ to $\tilde{p}$ by
a path $\tilde{\sigma}$, put $\pi(\tilde{\sigma})=\sigma$ and
$\beta=\sigma g\sigma^{-1}$, where, by abuse of notation, we also denote by $g$ a
path in the class $g$. In what follows, if $\beta:[0,1]\to M$ is a path in $M$,
$\tilde{\beta}_{\tilde{q}}(1)$ will denote the endpoint of the lifting to
$\tilde{M}$ of $\beta$ starting from $\tilde{q}$. Taking $\beta=\sigma g\sigma^{-1}$,
we define

$$
\alpha_{\tilde{p}}(\tilde{q})=\tilde{\beta}_{\tilde{q}}(1).
$$

\begin{definition}[free homotopy class]
  \label{def:dc-ch12-2-1}
  \dcref{ch12:2.1}
  A set $\mathcal{L}$ of closed paths in $M$ is called a \emph{free homotopy class} if
  given $f\in\mathcal{L}$ and $g:I\to M$ such that there exists a homotopy
  
  $$
  F:I\times I\to M,\quad F(0,t)=f(t),\ F(1,t)=g(t),\ F(s,0)=F(s,1),
  $$
  
  then $g\in\mathcal{L}$. The set of such classes will be denoted by $C_1(M)$.
  
  The difference between the definition above and the definition of the fundamental
  group is that in the free class we allow the origins of the paths to vary in $M$.
  
  The theorem below shows that in a compact Riemannian manifold $M$ with
  $\pi_1(M)\neq\{e\}$ there always exists a \emph{closed geodesic}, that is, a closed
  curve that is geodesic at all of its points. We should distinguish a closed
  geodesic from a \emph{geodesic lasso} which is a closed curve that is geodesic at all
  except one of its points, where it fails to be regular.
\end{definition}

\begin{theorem}[Cartan]
  \label{thm:dc-ch12-2-2}
  \dcref{ch12:2.2}
  If $M$ is compact and $\mathcal{L}\in C_1(M)$ is not the constant class, then
  there exists a closed geodesic of $M$ in the class $\mathcal{L}$.
\end{theorem}
\begin{proof}
  Let $d$ be the infimum of the lengths of piecewise differentiable curves
  belonging to $\mathcal{L}$. Since $\mathcal{L}$ is not trivial, $d>0$. Let
  $\gamma_j$ be a sequence of piecewise differentiable curves belonging to
  $\mathcal{L}$ such that $\ell(\gamma_j)\to d$. We can suppose that $\gamma_j$ is a
  broken geodesic defined on $[0,1]$ parametrized proportionally to arc length. Let
  $L=\sup\ell(\gamma_j)$. Then
  
  $$
  d(\gamma_j(t_1),\gamma_j(t_2))\leq\int_{t_1}^{t_2}|\gamma_j'(t)|\,dt\leq L(t_2-t_1),
  $$
  
  for all $t_1\leq t_2\in[0,1]$. Therefore the set $\{\gamma_j\}$ is equicontinuous.
  Since $M$ is compact, there exists a subsequence of $\gamma_j$, which we denote
  again by $\gamma_j$, which converges uniformly to a continuous closed curve
  $\gamma_o:[0,1]\to M$.
  
  Now let $0=t_o<t_1<\cdots<t_k=1$ be a partition of $[0,1]$ such that
  $\gamma_o|_{[t_{i-1},t_i]}$, $i=1,\dots,k$, is contained in a totally normal
  neighborhood. Let $\gamma:[0,1]\to M$ be a piecewise differentiable curve such
  that $\gamma^i=\gamma|_{[t_{i-1},t_i]}$ is the unique geodesic segment which joins
  the points $\gamma_o(t_{i-1})$ and $\gamma_o(t_i)$. It is clear that
  $\gamma\in\mathcal{L}$, hence $\ell(\gamma)\geq d$. We show that $\ell(\gamma)=d$.
  
  Suppose that $\ell(\gamma)>d$ and let $\varepsilon=\frac{\ell(\gamma)-d}{2k+1}$.
  There exists an integer $j$ such that
  
  $$
  \ell(\gamma_j)-d<\varepsilon\quad\text{and}\quad d(\gamma_j(t),\gamma_o(t))<\varepsilon,\ \text{for all }t\in[0,1].
  $$
  
  Denoting by $\gamma_j^i=\gamma_j|_{[t_{i-1},t_i]}$, we have
  
  $$
  \sum_{i=1}^k(\ell(\gamma_j^i)+2\varepsilon)=\ell(\gamma_j)+2k\varepsilon<d+(2k+1)\varepsilon=\ell(\gamma)=\sum_{i=1}^k\ell(\gamma^i).
  $$
  
  Therefore, there exists an integer $i$, $1\leq i\leq k$, such that
  $\ell(\gamma_j^i)+2\varepsilon<\ell(\gamma^i)$, which contradicts the fact that
  $\gamma^i$ is minimizing and proves that $\ell(\gamma)=d$.
  
  We parametrize $\gamma$ by arc length. Then $\gamma:[0,d]\to M$ is a broken
  geodesic which has minimum length in the class $\mathcal{L}$. We show that
  $\gamma$ is regular at the point $p_i=\gamma(t_i)$, for all $i=0,\dots,k$. Suppose
  to the contrary and let $B$ be a convex ball centered at $p_i$. Choose points
  $q_1$ and $q_2$ in $\gamma\cap B$ in a way that the geodesic triangle
  $p_iq_1q_2$ is homotopic to a point. Then the closed curve constituted by the
  minimizing geodesic $q_1q_2$ and by the arc of $\gamma$ between $q_1$ and $q_2$
  that does not contain $p_i$ is in the class $\mathcal{L}$ and has length smaller
  than $\gamma$, which is a contradiction. $\blacksquare$
\end{proof}

\begin{remark}[Remark 2.3]
  \label{rem:dc-ch12-2-3}
  \dcref{ch12:2.3}
  If $M$ is not compact, the theorem is false, as shown by the example of a surface
  of revolution generated by a curve which is asymptotic to the axis of revolution.
  Since there exist curves of arbitrarily small length in the free homotopy classes
  which are non-trivial, such classes do not contain curves of minimum length.
\end{remark}

\begin{remark}[Remark 2.4]
  \label{rem:dc-ch12-2-4}
  \dcref{ch12:2.4}
  If $M$ is simply connected (and compact), the existence of a closed geodesic in
  $M$, although true, is a more difficult problem. Indeed, the problem of
  determining the number and the nature of the closed geodesics on a Riemannian
  manifold is one of the most celebrated chapters in Geometry. A good reference on
  this topic is Klingenberg [Kℓ 4].
  
  For our purposes, it is important to know when an isometry of $\tilde{M}$ without
  fixed points leaves invariant a geodesic.
\end{remark}

\begin{definition}[translation]
  \label{def:dc-ch12-2-5}
  \dcref{ch12:2.5}
  An isometry $f:\tilde{M}\to\tilde{M}$ without fixed points is said to be a
  \emph{translation} of $\tilde{M}$ if it leaves invariant some geodesic $\tilde{\gamma}$
  of $\tilde{M}$, that is, if $f(c)=c$, where $c=\tilde{\gamma}((-\infty,\infty))$.
  In this case, we say that $f$ is a \emph{translation along} $\tilde{\gamma}$.
\end{definition}

\begin{proposition}[covering transformations are translations]
  \label{prop:dc-ch12-2-6}
  \dcref{ch12:2.6}
  Let $M$ be a compact Riemannian manifold and let $\alpha\neq\mathrm{ident}$ be a covering transformation
  of $\tilde{M}$, considered with the covering metric. Then $\alpha$ is a
  translation of $\tilde{M}$.
\end{proposition}
\begin{proof}
  Let $\tilde{p}\in\tilde{M}$ and let $g\in\pi_1(M;p)$, $p=\pi(\tilde{p})$,
  the element corresponding to $\alpha$ under the isomorphism mentioned in the
  introduction of this Section. By
  Cartan's Theorem (2.2), there exists a closed geodesic $\gamma$ of $M$ in the
  free homotopy class determined by $g$. Choose a point $q\in\gamma$. Then $\gamma$
  is homotopic to the closed path $\sigma g\sigma^{-1}$, where $\sigma$ is a path
  joining $p$ to $q$. Let $\tilde{q}=\tilde{\sigma}_{\tilde{p}}(1)$, that is,
  $\tilde{q}$ is the endpoint of the lifting of $\sigma$ starting from $\tilde{p}$.
  Let $\tilde{\gamma}$ be the lift of $\gamma$ starting from $\tilde{q}$; we show
  that $\alpha$ leaves $\tilde{\gamma}$ invariant.
  
  For this, let $\alpha_{\tilde{q}}\in A(\tilde{M})$ be the isometry corresponding
  to the class $[\gamma]\in\pi_1(M;q)$ and the point $\tilde{q}$, in the isomorphism
  indicated above. We claim that $\alpha_{\tilde{q}}=\alpha$. In fact, since
  $\gamma$ is homotopic to $\sigma g\sigma^{-1}$, its lift starting from $\tilde{q}$
  has the same endpoint, that is,
  
  $$
  \alpha(\tilde{q})=\alpha_{\tilde{q}}(\tilde{q}).
  $$
  
  Therefore, $\tilde{q}$ is a fixed point of $\alpha\circ\alpha_{\tilde{q}}^{-1}$,
  and the claim follows easily.
  
  It follows that if $\tilde{\gamma}(s)$ is a point of the lift of $\gamma$ starting
  from $\tilde{q}$, we have, by uniqueness of lifting,
  
  $$
  \alpha(\tilde{\gamma}(s))=\alpha_{\tilde{q}}(\tilde{\gamma}(s))\in\tilde{\gamma},
  $$
  
  which shows that $\tilde{\gamma}$ is invariant by $\alpha$, and proves the
  Proposition. $\blacksquare$
\end{proof}

\section{Preissman's Theorem}

A \emph{geodesic triangle} $T$ in a Riemannian manifold $M$ is a set formed by three
segments of minimizing normalized geodesics (called \emph{sides} of the triangle)

$$
\gamma_1:[0,\ell_1]\to M,\quad\gamma_2:[0,\ell_2]\to M,\quad\gamma_3:[0,\ell_3]\to M,
$$

in such a way that $\gamma_i(\ell_i)=\gamma_{i+1}(0)$, $i=1,2$ and
$\gamma_3(\ell_3)=\gamma_1(0)$. The endpoints of the geodesic segments are called
\emph{vertices} of $T$. The angle

$$
\sphericalangle(-\gamma_i'(\ell_i),\gamma_{i+1}'(0)),\ i=1,2,\quad\text{or}\quad\sphericalangle(-\gamma_3'(\ell_3),\gamma_1'(0)),
$$

is called the (interior) \emph{angle} of the corresponding vertex.

\begin{lemma}[comparison for geodesic triangles, $K\leq 0$]
  \label{lem:dc-ch12-3-1}
  \dcref{ch12:3.1}
  \uses{prop:dc-ch10-2-5}
  Let $\tilde{M}$ be a complete simply connected Riemannian manifold, with curvature
  $K\leq 0$. Let $a$, $b$ and $c$ be three points of $\tilde{M}$. Such points
  determine a unique geodesic triangle $T$ in $\tilde{M}$ with vertices $a$, $b$,
  $c$. Let $\alpha$, $\beta$ and $\gamma$ be the angles of the vertices $a$, $b$,
  $c$, respectively, and let $A$, $B$, $C$ be the lengths of the sides opposite the
  vertices $a$, $b$, $c$, respectively. Then
  - (i) $A^2+B^2-2AB\cos\gamma\leq C^2$ ($<C^2$, if $K<0$).
  - (ii) $\alpha+\beta+\gamma\leq\pi$ ($<\pi$, if $K<0$).
\end{lemma}
\begin{proof}
  Let $\gamma_A$, $\gamma_B$, $\gamma_C$ be the geodesics of lengths
  $\ell(\gamma_A)=A$, $\ell(\gamma_B)=B$, $\ell(\gamma_C)=C$ which form the sides of
  $T$. Let $\Gamma_A=\exp_c^{-1}(\gamma_A)$, $\Gamma_B=\exp_c^{-1}(\gamma_B)$ and
  $\Gamma_C=\exp_c^{-1}(\gamma_C)$ be curves in $T_c(\tilde{M})$. Since $\gamma_A$
  and $\gamma_B$ are radial geodesics issuing from the origin $c$, we have
  
  $$
  A=\ell(\gamma_A)=\ell(\Gamma_A),\qquad B=\ell(\gamma_B)=\ell(\Gamma_B).
  $$
  
  In addition, denoting by $\Gamma_o$ the segment of the straight line in
  $T_c(\tilde{M})$ that joins the endpoints of $\Gamma_C$, we have that
  $\ell(\Gamma_o)\leq\ell(\Gamma_C)$ and
  
  $$
  \ell(\Gamma_o)^2=A^2+B^2-2AB\cos\gamma.
  $$
  
  Since $K\leq 0$ and $T_c(\tilde{M})$ has zero curvature, we can apply \cref{prop:dc-ch10-2-5} (application of Rauch's Theorem) and obtain that
  
  $$
  \ell(\Gamma_C)\leq\ell(\gamma_C),
  $$
  
  with strict inequality if $K<0$.
  It follows that
  
  $$
  A^2+B^2-2AB\cos\gamma\leq\ell(\Gamma_C)^2\leq\ell(\gamma_C)^2=C^2,
  $$
  
  where the last inequality is strict if $K<0$.
  This proves (i).
  
  To prove (ii), let us observe that
  
  $$
  C=d(a,b),\qquad B=d(a,c),\qquad A=d(b,c),
  $$
  
  and, therefore, each length $A$, $B$ or $C$ is bounded by the sum of the other
  two. We can then consider in the Euclidean space $T_c(\tilde{M})$ a triangle whose
  sides have lengths $A$, $B$ and $C$. Denoting the opposite angles of this triangle
  by $\alpha'$, $\beta'$ and $\gamma'$, respectively, we obtain from (i),
  
  $$
  \alpha\leq\alpha',\qquad\beta\leq\beta',\qquad\gamma\leq\gamma'.
  $$
  
  All three inequalities are strict if $K<0$.
  Since $\alpha'+\beta'+\gamma'=\pi$, (ii) follows. $\blacksquare$
  
  From now on, $M$ will denote a complete Riemannian manifold with sectional
  curvature $K<0$. As always, $\pi:\tilde{M}\to M$ denotes the universal covering of
  $M$ with the covering metric. Our goal is to prove the following theorem.
\end{proof}

\begin{theorem}[Preissman]
  \label{thm:dc-ch12-3-2}
  \dcref{ch12:3.2}
  If $M$ is a compact Riemannian manifold with negative curvature, then any abelian
  subgroup of the fundamental group $\pi_1(M)$, different from the identity, is
  infinite cyclic.
  
  The proof depends on a sequence of lemmas.
\end{theorem}

\begin{lemma}[uniqueness of the invariant geodesic]
  \label{lem:dc-ch12-3-3}
  \dcref{ch12:3.3}
  \uses{lem:dc-ch12-3-1}
  If $K<0$ and $f:\tilde{M}\to\tilde{M}$ is a translation along the geodesic
  $\tilde{\gamma}$, $f\neq\mathrm{id}$, then $\tilde{\gamma}$ is the unique geodesic
  left invariant by $f$.
\end{lemma}
\begin{proof}
  Suppose that $f$ leaves invariant two geodesics $\tilde{\gamma}_1$ and
  $\tilde{\gamma}_2$. Since $f$ does not have fixed points,
  $\tilde{\gamma}_1\cap\tilde{\gamma}_2=\emptyset$; otherwise,
  $\tilde{\gamma}_1\cap\tilde{\gamma}_2$ has at least two points, which contradicts
  the simple connectivity of $\tilde{M}$. Let $\tilde{p}_1\in\tilde{\gamma}_1$,
  $\tilde{p}_2\in\tilde{\gamma}_2$, and let $\tilde{\gamma}_3$ be the minimizing
  geodesic joining $\tilde{p}_1$ to $\tilde{p}_2$. Consider the "geodesic
  quadrilateral" $\tilde{p}_1$, $f(\tilde{p}_1)$, $\tilde{p}_2$, $f(\tilde{p}_2)$ in
  $\tilde{M}$ and denote its (interior) angles by $\alpha$, $\pi-\alpha'$ (adjacent
  to the side $\tilde{\gamma}_1$), $\beta$, $\pi-\beta'$ (adjacent to the side
  $\tilde{\gamma}_2$). Since $f$ is an isometry, $\alpha=\alpha'$ and $\beta=\beta'$.
  Therefore, the sum of the interior angles of such a quadrilateral is equal to
  $2\pi$. Now divide this quadrilateral into two triangles $T_1$ and $T_2$ and
  denote by $\sum_i$ the sum of the interior angles of the triangle $T_i$,
  $i=1,2$. It is clear that, at each vertex, the sum of the angles of the triangles
  which meet there is greater than or equal to the angle of the quadrilateral at
  this vertex. Therefore $\sum_1+\sum_2\geq 2\pi$. It follows that one of the two
  triangles has its sum of interior angles $\geq\pi$, which contradicts \cref{lem:dc-ch12-3-1}
  (ii). $\blacksquare$
\end{proof}

\begin{lemma}[commuting isometry is a translation along the same geodesic]
  \label{lem:dc-ch12-3-4}
  \dcref{ch12:3.4}
  If $K<0$ and $g:\tilde{M}\to\tilde{M}$ is an isometry without fixed points which
  commutes with a translation $f$ along $\tilde{\gamma}$, $f\neq\mathrm{id}$, then
  $g$ is a translation along $\tilde{\gamma}$.
\end{lemma}
\begin{proof}
  It suffices to observe that
  
  $$
  f\circ g(\tilde{\gamma})=g\circ f(\tilde{\gamma})=g(\tilde{\gamma}),
  $$
  
  hence, by uniqueness of the previous lemma (3.3), $g(\tilde{\gamma})=\tilde{\gamma}$.
  $\blacksquare$
\end{proof}

\begin{lemma}[subgroup leaving a geodesic invariant is infinite cyclic]
  \label{lem:dc-ch12-3-5}
  \dcref{ch12:3.5}
  \uses{prop:dc-ch12-2-6}
  If all the elements of a non-trivial subgroup $H\subset\pi_1(M)$, considered as
  isometries of $\tilde{M}$, leave invariant a fixed geodesic $\tilde{\gamma}$, then
  $H$ is infinite cyclic.
\end{lemma}
\begin{proof}
  Fix a point $\tilde{p}\in\tilde{\gamma}$ as origin and consider the map
  $\theta:H\to\mathbb{R}$ given by $\theta(h)=\pm d(\tilde{p},h(\tilde{p}))$, where
  the sign $-$ or $+$ is used accordingly as $h(\tilde{p})$ is "before" or "after"
  $\tilde{p}$, respectively, in the orientation of $\tilde{\gamma}$. Since the
  elements of $H$ are isometries which leave $\tilde{\gamma}$ invariant, $\theta$ is
  a homomorphism of $H$ into the additive group of the reals $\mathbb{R}$. $\theta$
  is an injective homomorphism; otherwise, $h_1(\tilde{p})=h_2(\tilde{p})$, with
  $h_1\neq h_2$, hence $\tilde{p}$ would be a fixed point of $h_1h_2^{-1}$.
  Therefore, $H$ is an additive subgroup of $\mathbb{R}$. But any additive subgroup
  of $\mathbb{R}$ is either dense in $\mathbb{R}$ or infinite cyclic. Since the
  isometries of $H$ operate in a totally discontinuous way, $H$ is not dense. Hence
  $H$ is infinite cyclic. $\blacksquare$
  
  \emph{Proof of Preissman's Theorem (3.2).} Let $H\subset\pi_1(M)$ be an abelian
  subgroup with $H\neq\{e\}$. By \cref{prop:dc-ch12-2-6} and by Lemmas 3.3 and 3.4, there
  exists a unique geodesic $\tilde{\gamma}$ which is invariant under all the
  elements of $H$. By Lemma 3.5, $H$ is infinite cyclic. $\blacksquare$
\end{proof}

\begin{example}[surface of genus two times a circle]
  \label{ex:dc-ch12-3-6}
  \dcref{ch12:3.6}
  Let $N$ be a surface of genus two and let $M=N\times S^1$ be the product manifold
  of $N$ with a circle $S^1$. Then $M$ does not carry a metric of negative
  curvature. Indeed, if $C\subset\pi_1(N)$ is a cyclic subgroup, then
  $C\oplus\mathbb{Z}\subset\pi_1(M)$ is an abelian subgroup which is not cyclic.
\end{example}

\begin{example}[$m$-torus]
  \label{ex:dc-ch12-3-7}
  \dcref{ch12:3.7}
  The $m$-torus $S^1\times\cdots\times S^1=T^m$, $m\geq 2$, does not carry a metric
  of negative curvature, since its fundamental group is
  $\mathbb{Z}\oplus\cdots\oplus\mathbb{Z}$ ($m$ times).
  
  The same technique used in Preissman's Theorem allows us to obtain a slightly more
  general theorem. To that end, we need another theorem on the fundamental group,
  also due to Preissman.
\end{example}

\begin{theorem}[$\pi_1$ not abelian]
  \label{thm:dc-ch12-3-8}
  \dcref{ch12:3.8}
  If $M$ is compact and $K<0$, then $\pi_1(M)$ is not abelian.
\end{theorem}
\begin{proof}
  To prove what is needed, we establish the lemma below which is a little
  more general.
\end{proof}

\begin{lemma}[invariant geodesic forces non-compactness]
  \label{lem:dc-ch12-3-9}
  \dcref{ch12:3.9}
  \uses{lem:dc-ch12-3-1}
  If $M$ is complete, $K\leq 0$, and there exists a geodesic invariant under the
  elements of $A(\tilde{M})$, then $M$ is not compact.
  
  \emph{Proof of the Lemma.} Let $\tilde{\gamma}$ be the geodesic invariant under the
  elements of $A(\tilde{M})$. Fix a point $\tilde{p}\in\tilde{\gamma}$, a real
  number $t>0$, and consider the normalized geodesic $\tilde{\beta}:[0,t]\to\tilde{M}$,
  $\tilde{\beta}(0)=\tilde{p}$, perpendicular to $\tilde{\gamma}$ at $\tilde{p}$.
  Let $\beta=\pi\circ\tilde{\beta}$, $\gamma=\pi\circ\tilde{\gamma}$,
  $p=\pi(\tilde{p})$, and $\alpha_t$ be a minimizing geodesic in $M$ joining
  $\beta(t)$ to $p$. We show that the length $\ell(\alpha_t)=t$.
  
  Let $\tilde{\alpha}_t$ be the lift of $\alpha_t$ starting from $\tilde{\beta}(t)$.
  Since $\tilde{\gamma}$ is invariant, the endpoint of $\tilde{\alpha}_t$ belongs to
  $\tilde{\gamma}$. Since $K\leq 0$, by \cref{lem:dc-ch12-3-1} (i),
  $\ell(\tilde{\alpha}_t)\geq\ell(\tilde{\beta})$. On the other hand,
  
  $$
  \ell(\tilde{\alpha}_t)=\ell(\alpha_t)\leq\ell(\beta)=\ell(\tilde{\beta})=t,
  $$
  
  hence $\ell(\alpha_t)=t$, as was claimed.
  
  Since $t$ is arbitrary, $M$ is not bounded, which concludes the proof of Lemma
  3.9. $\blacksquare$
  
  The theorem (3.8) follows now from the fact that if $M$ is compact, $K<0$ and
  $\pi_1(M)$ is abelian, there exists a geodesic invariant under all the elements of
  $A(\tilde{M})$. $\blacksquare$
  
  We can now refine Preissman's Theorem and prove the following fact.
\end{lemma}

\begin{theorem}[Byers]
  \label{thm:dc-ch12-3-10}
  \dcref{ch12:3.10}
  \uses{thm:dc-ch12-3-8, lem:dc-ch12-3-5}
  If $M$ is compact, $K<0$, and $H$ is a solvable subgroup of $\pi_1(M)$,
  $H\neq\{e\}$, then $H$ is infinite cyclic. In addition, $\pi_1(M)$ does not have a
  cyclic subgroup of finite index.
\end{theorem}
\begin{proof}
  Since $H$ is solvable, choose a finite sequence of subgroups with $k$
  minimal:
  
  $$
  H=H_0\supset H_1\supset\cdots\supset H_{k-1}\supset H_k=\{e\}
  $$
  
  such that $H_{i+1}$ is normal in $H_i$ and $H_i/H_{i+1}$ is abelian. Then
  $H_{k-1}\neq\{e\}$ and $H_{k-1}$ is abelian, hence, by Preissman's Theorem (3.2),
  infinite cyclic. If $k=1$, this proves the first assertion. Assume $k\geq2$.
  
  Let $g\in\pi_1(M)$ be a generator of $H_{k-1}$ and let $\tilde{\gamma}$ be the
  geodesic in $\tilde{M}$ invariant under $g$. Let $a\in H_{k-2}$, and choose a
  nonidentity element $b\in H_{k-1}$, for instance $b=g$.
  Since $a^{-1}b^{-1}ab\in H_{k-1}$, we have, for some integer $n$,
  
  $$
  a^{-1}b^{-1}ab=g^n.
  $$
  
  It follows that
  
  $$
  a^{-1}b^{-1}ab(\tilde{\gamma})=g^n(\tilde{\gamma})=\tilde{\gamma}.
  $$
  
  Since $b\in H_{k-1}$, $b(\tilde{\gamma})=\tilde{\gamma}$. Therefore
  $b^{-1}a(\tilde{\gamma})=a(\tilde{\gamma})$, that is, $b^{-1}$ leaves invariant the
  geodesic $a(\tilde{\gamma})$. By uniqueness, $a(\tilde{\gamma})=\tilde{\gamma}$.
  Therefore, all the elements of $H_{k-2}$ leave $\tilde{\gamma}$ invariant. By
  \cref{lem:dc-ch12-3-5}, $H_{k-2}$ is infinite cyclic.
  
  Repeating the argument above a finite number of times, we conclude that $H$ is
  infinite cyclic, which proves the first statement of the theorem.
  
  To prove the final assertion, suppose that there exists a subgroup
  $H\subset\pi_1(M)$, cyclic and of finite index. If $H=\{e\}$, then $\pi_1(M)$ is
  finite; this contradicts Preissman's Theorem (3.2) unless $\pi_1(M)=\{e\}$, which
  contradicts \cref{thm:dc-ch12-3-8}. Thus $H=\langle g\rangle$ with $g\neq e$. If
  $H=\pi_1(M)$, then $\pi_1(M)$ is abelian, again contradicting \cref{thm:dc-ch12-3-8}. Hence
  $H$ is proper. Let $\tilde{\gamma}$ be the geodesic of $\tilde{M}$ invariant under
  $g$, and let $a\in\pi_1(M)-H$. Since $H$ has finite index, choose $n>0$ with
  $a^n\in H$. Applying Preissman's Theorem to the subgroup generated by $a$, we have
  $a^n\neq e$, so $a^n=g^m$ for some $m\neq0$. Therefore
  
  $$
  a^n(\tilde{\gamma})=g^m(\tilde{\gamma})=\tilde{\gamma}.
  $$
  
  The nonidentity transformation $a^n$ also leaves $a(\tilde{\gamma})$ invariant, so
  by uniqueness, $a(\tilde{\gamma})=\tilde{\gamma}$, for all $a\in\pi_1(M)-H$. It
  follows that every element of $\pi_1(M)$ leaves invariant the geodesic
  $\tilde{\gamma}$. By \cref{lem:dc-ch12-3-5}, $\pi_1(M)$ is infinite cyclic, which contradicts
  \cref{thm:dc-ch12-3-8}. $\blacksquare$
\end{proof}

\begin{remark}[Remark 3.11]
  \label{rem:dc-ch12-3-11}
  \dcref{ch12:3.11}
  Riemannian manifolds of non-positive curvature form a substantial topic in
  Riemannian Geometry, which we have barely touched. To the reader interested in
  some of the recent developments, we recommend the excellent survey by P. Eberlein,
  "Structure of manifolds of nonpositive curvature", in *Global Geometry and Global
  Analysis 1984*, Proceedings, Berlin, edited by D. Ferus, R. Gardner, S. Helgason
  and U. Simon, Lecture Notes in Math. 1156, Springer Verlag, 1985, pp. 86–153.
  Other relationships between the geometry of a manifold $M$ of
  non-positive curvature and $\pi_1(M)$ are described in Section 7 of this survey.
\end{remark}
